%% kapa: exam=1 round=37 year=2026 period=1 form=A subject=civil_law subject_ko=민법 seq=38 number=38
%% source: Q-Net 감정평가사 기출문제 (cst00309) artlSeq=5258924 file=raw/1008/5258924/2026년 제37회 감정평가사 1차 1교시.pdf
%% conversion: kordoc 4.24.51 -> tools/gen1.py (auto); review=proofread-sonnet; flags=none
\begin{question}{38}
乙이 甲으로부터 100만 원을 차용하면서 그가 착용하고 있던 X시계에 관하여 甲에게 질권을 설정한 경우에 관한 설명으로 옳지 않은 것은? (다툼이 있으면 판례에 따름)
\begin{choices}
  \item 甲은 선량한 관리자의 주의로 X시계를 점유하여야 한다.
  \item 乙에게 X시계에 관한 처분권한이 없었다는 점에 관하여 甲이 선의ㆍ무과실이었다면, 甲은 질권을 선의취득할 수 있다.
  \item 甲은 乙에 대한 채권의 일부를 변제받았더라도 X시계 전부에 대하여 그 권리를 행사할 수 있다.
  \item 甲은 乙에 대한 채권의 변제를 받을 때까지 X시계를 유치할 수 있으나, 자기보다 우선권이 있는 채권자에게 대항하지 못한다.
  \item 甲은 乙에 대한 채권의 변제를 받기 위하여 X시계를 경매할 수 있고, 그 매각대금으로부터 일반채권자와 동일한 순위로 변제받는다.
\end{choices}
\end{question}
\answer{5}
